\chapter{Gynaecomastia Surgery}

\section*{Introduction \& Clinical Indications}
Gynaecomastia surgery is a reconstructive plastic surgery procedure designed to reduce enlarged male breast tissue that has not resolved spontaneously or through medical management. This condition, characterized by an excess of glandular tissue, fatty tissue, or both, can affect one or both breasts and often leads to significant physical discomfort, self-consciousness, and psychological distress. The surgical intervention typically involves a combination of liposuction to remove excess adipose tissue and direct surgical excision of the firm, fibrous glandular tissue that is resistant to liposuction. The primary objective is to achieve a flatter, firmer, and more masculine chest contour that is proportionate to the individual's body habitus, thereby improving both aesthetic appearance and quality of life.

\begin{itemize}
  \item Male breast reduction
  \item Gynecomastia excision
  \item Glandular tissue excision with liposuction
  \item Subcutaneous mastectomy for gynecomastia
  \item Chest masculinization surgery (in some contexts)
  \item Mammaplasty, reduction, male
\end{itemize}

The underlying principle of gynaecomastia surgery is the selective removal of excess breast tissue components to restore a natural male chest contour. True gynaecomastia involves hyperplasia of the glandular tissue, which is firm, often retroareolar, and cannot be effectively removed by liposuction alone. Pseudogynaecomastia, conversely, is primarily composed of excess fatty tissue. Most patients present with a mixed picture, necessitating a dual approach. Liposuction, often performed using a tumescent technique, targets the fatty component, allowing for precise contouring and removal of diffuse adipose tissue through small, inconspicuous incisions. Concurrently, direct surgical excision, typically via a periareolar incision, is employed to remove the dense, subareolar glandular disc and any redundant skin.

This combined strategy ensures comprehensive reduction of both fatty and glandular elements, minimizes scarring, and facilitates optimal skin redraping. The rationale for this approach is to achieve a smooth, symmetric, and aesthetically pleasing result by addressing all components of the enlargement. Technical goals include achieving symmetry between the two sides, eliminating breast mound projection, minimizing nipple-areola complex distortion, and creating a well-defined pectoral contour that is proportionate to the patient's physique. The choice of technique is tailored to the individual's specific anatomy, including the amount of glandular tissue, fat, and skin laxity.

The treatment of gynaecomastia has evolved significantly over centuries. Early attempts to address enlarged male breasts date back to ancient civilizations, with rudimentary excisional techniques described in historical medical texts. In the late 19th and early 20th centuries, surgeons like William Halsted and Willy Meyer described more systematic approaches involving direct excision, often resulting in significant scarring due to the removal of large skin segments. These early techniques prioritized tissue removal over aesthetic outcome.

The modern era of gynaecomastia surgery began to take shape with the advent of liposuction in the 1980s. Pioneered by surgeons such as Yves-G\'{e}rard Illouz and Pierre Fournier, liposuction revolutionized body contouring, and its application to gynaecomastia allowed for less invasive fat removal. Subsequent refinements, including the development of the tumescent technique by Jeffrey Klein in the 1980s and the introduction of ultrasound-assisted liposuction (UAL) and power-assisted liposuction (PAL) in the 1990s, further improved the precision of fat removal and reduced operative trauma. Contemporary techniques emphasize minimal scarring, often utilizing periareolar incisions for glandular excision combined with small stab incisions for liposuction, aiming for superior aesthetic outcomes and faster recovery. The evolution reflects a shift from purely excisional, often disfiguring, procedures to sophisticated, minimally invasive techniques that prioritize both effective tissue reduction and aesthetic contouring.

Surgery for gynaecomastia is typically indicated for specific clinical scenarios after a thorough diagnostic workup has ruled out other causes and non-surgical options have been exhausted.

\begin{itemize}
  \item Persistent, symptomatic gynaecomastia that has not resolved spontaneously, particularly in adolescents after 18 months of observation or in adults.
  \item Significant physical discomfort, including breast tenderness, pain, or nipple sensitivity, especially with clothing friction or during physical activity.
  \item Psychological distress, embarrassment, or social anxiety due to the appearance of enlarged breasts, severely impacting self-esteem, body image, and social interactions.
  \item Inability to participate in certain physical activities or wear desired clothing due to the breast enlargement.
  \item Failure of medical management, such as discontinuation of causative medications or treatment of underlying hormonal imbalances, to reduce breast size.
  \item Presence of firm, palpable glandular tissue that is unlikely to respond to non-surgical interventions, often confirmed by imaging.
  \item Significant asymmetry of the breast tissue that causes aesthetic concern and functional imbalance.
  \item Suspicion of malignancy, although surgery for gynaecomastia is primarily for benign conditions, any suspicious mass warrants biopsy or excision.
\end{itemize}

Gynaecomastia surgery is considered a primary surgical treatment for persistent and symptomatic breast enlargement in males. However, in the broader clinical management pathway, it functions as a secondary intervention. This is because a comprehensive medical evaluation is always the primary step to identify and address any underlying causes of gynaecomastia. Such causes can include hormonal imbalances (e.g., testosterone deficiency, estrogen excess), certain medications (e.g., spironolactone, cimetidine, some antidepressants), illicit drug use (e.g., marijuana, anabolic steroids), liver or kidney disease, and rarely, tumors.

If an underlying medical condition or medication is identified, the primary approach involves treating the condition or discontinuing the offending drug. Only after these non-surgical measures have been exhausted or deemed inappropriate, or if the gynaecomastia is idiopathic and persistent, does surgery become the definitive primary treatment option. For adolescents, a period of observation (typically 12-18 months) is often recommended as pubertal gynaecomastia frequently resolves spontaneously. Surgery is positioned as a definitive solution when the condition is stable, bothersome, and unlikely to regress further.

\section*{Relevant Surgical Anatomy}
The male breast, though typically less developed than the female breast, shares fundamental anatomical structures crucial for gynaecomastia surgery. It consists primarily of glandular tissue, adipose tissue, skin, and the nipple-areola complex (NAC), all overlying the pectoralis major muscle. The glandular tissue, which proliferates in true gynaecomastia, is typically concentrated in a retroareolar disc, firm and fibrous. Adipose tissue surrounds this glandular component and extends into the chest wall. The skin of the male chest is generally thicker than that of the female breast, with less elasticity, which influences redraping post-excision. The nipple-areola complex is a key aesthetic landmark, and its position, size, and projection are critical considerations during surgery.

The arterial supply to the male breast is derived predominantly from perforating branches of the internal mammary artery (medially), the lateral thoracic artery (laterally), and intercostal perforators. Venous drainage largely parallels the arterial supply, emptying into the internal mammary and axillary veins. Innervation of the chest wall and nipple-areola complex is primarily from the anterior and lateral cutaneous branches of the intercostal nerves (T2-T6). Preservation of these nerves is vital for maintaining nipple sensation. Lymphatic drainage follows the major blood vessels, primarily to the axillary lymph nodes, with some drainage to the internal mammary nodes. Surgical landmarks include the inframammary fold, the sternal notch, and the mid-clavicular line, which guide incision placement and contouring to achieve a natural male chest aesthetic.

\section*{Preoperative Workup \& Preparation}
Thorough preoperative preparation is crucial for optimizing surgical outcomes and patient safety in gynaecomastia surgery.

\begin{itemize}
  \item \textbf{Medical Evaluation:} A comprehensive medical history and physical examination are performed to assess overall health. This includes a detailed review of medications, supplements, and any history of illicit drug use (especially anabolic steroids).
  \item \textbf{Laboratory Tests:} Blood tests are routinely performed to assess general health, including a complete blood count, basic metabolic panel, and coagulation studies. Hormone levels (e.g., testosterone, estrogen, LH, FSH, prolactin, thyroid function) are often checked to rule out endocrine causes.
  \item \textbf{Imaging Studies:} Mammography or ultrasound of the breast may be performed, especially in older patients or those with suspicious findings (e.g., unilateral, firm, fixed mass), to differentiate benign gynaecomastia from breast cancer.
  \item \textbf{Anesthesia Clearance:} Patients typically undergo a preoperative assessment by an anesthesiologist to ensure they are fit for general anesthesia. This may include an electrocardiogram (ECG) and chest X-ray for certain age groups or medical conditions.
  \item \textbf{Medication Adjustments:} Patients are instructed to stop smoking several weeks before surgery to improve wound healing and reduce complication risks. Anti-inflammatory drugs (NSAIDs), aspirin, herbal supplements (e.g., ginkgo, ginseng), and other blood-thinning medications should be discontinued at least two weeks prior to surgery to reduce bleeding risk.
  \item \textbf{NPO Status:} Patients must adhere to NPO (nil per os) guidelines, typically refraining from food and drink for 6-8 hours before surgery to prevent aspiration during anesthesia.
  \item \textbf{Prophylactic Antibiotics:} A single dose of intravenous prophylactic antibiotics is typically administered preoperatively to reduce the risk of surgical site infection.
  \item \textbf{Informed Consent:} A detailed discussion of the procedure, potential risks, expected benefits, and alternative treatments is conducted, followed by obtaining informed consent from the patient.
  \item \textbf{Logistics:} Patients are advised to arrange for a responsible adult to drive them home after surgery and assist with initial recovery. Planning for adequate time off from work or school is also important.
\end{itemize}

Careful patient selection is paramount to ensure safety and satisfactory results in gynaecomastia surgery.

\textbf{Situations requiring diagnostic clarification, optimization, or an alternative plan:}
\begin{itemize}
  \item Active breast infection or inflammation, which must be treated before surgery.
  \item Uncontrolled systemic medical conditions, such as severe heart disease, uncontrolled diabetes, severe coagulopathy, or significant respiratory compromise.
  \item Untreated or undiagnosed breast malignancy; any suspicious mass requires definitive diagnosis prior to elective surgery.
  \item Unrealistic patient expectations regarding surgical outcomes, which can lead to severe dissatisfaction.
  \item Active illicit drug use, particularly anabolic steroids, which can cause recurrence if not ceased and may pose anesthetic risks.
\end{itemize}

\textbf{Relative Contraindications \& Precautions:}
\begin{itemize}
  \item Significant obesity: While not an absolute contraindication, extreme obesity can complicate surgery, increase risks, and compromise aesthetic results. Weight loss is often recommended preoperatively.
  \item Ongoing use of medications known to cause gynaecomastia, if they cannot be safely discontinued or replaced.
  \item Unstable psychological conditions or body dysmorphic disorder, which may lead to dissatisfaction despite a technically successful surgery. Psychiatric clearance may be required.
  \item Adolescents with pubertal gynaecomastia that is still actively regressing; observation for 12-18 months is often preferred to allow for spontaneous resolution.
  \item History of keloid or hypertrophic scarring, requiring careful incision placement, meticulous closure, and aggressive postoperative scar management.
  \item Smoking: Significantly increases risks of poor wound healing, skin necrosis, infection, and other complications. Patients are strongly advised to quit several weeks prior to surgery.
  \item Bleeding disorders or use of anticoagulants, requiring careful management and potential temporary cessation of medication.
\end{itemize}

\section*{Surgical Technique \& Procedure}
Gynaecomastia surgery is typically performed under general anesthesia, though local anesthesia with sedation may be an option for very minor cases. The duration usually ranges from one to three hours, depending on the extent of tissue removal and the complexity of the case.

\begin{enumerate}
  \item  \textbf{Anesthesia and Patient Positioning:} The patient is placed in a supine position with arms abducted to 90 degrees on arm boards. General anesthesia is most common, ensuring patient comfort and immobility throughout the procedure. The chest area, including the axillae and abdomen, is prepped with an antiseptic solution and draped in a sterile fashion. Preoperative markings are crucial, outlining the extent of glandular tissue, fatty deposits, and the desired chest contour.

  2.  \textbf{Tumescent Infiltration:} A tumescent solution, typically a mixture of normal saline, lidocaine (for local anesthesia), and epinephrine (for vasoconstriction), is infiltrated into the subcutaneous fatty tissue of the chest. This solution numbs the area, constricts blood vessels to minimize bleeding, and hydro-dissects the fat, making it easier to remove via liposuction.

  3.  \textbf{Liposuction:} Small incisions, typically 3-5 mm in length, are made in inconspicuous areas such as the lateral chest wall, axillary fold, or inframammary fold. Through these incisions, thin liposuction cannulas are inserted. The surgeon meticulously suctions out excess adipose tissue, sculpting the chest contour to achieve a smooth transition from the chest to the abdomen and axilla, and reducing the overall volume.

  4.  \textbf{Glandular Excision:} A periareolar incision, most commonly a Webster incision at the inferior border of the areola, is made. This incision provides direct access to the firm, subareolar glandular tissue. The glandular disc is carefully dissected from the overlying skin and the underlying pectoralis major muscle fascia using electrocautery or a scalpel. The entire glandular component is removed, ensuring a flat, firm base.

  5.  \textbf{Skin Redraping and Nipple-Areola Complex Management:} For significant gynaecomastia with skin laxity, additional skin excision may be necessary. This can involve concentric periareolar excisions or, in severe cases, a free nipple graft technique. The remaining skin is then redraped to achieve a smooth contour. The nipple-areola complex may be resized or repositioned if indicated to ensure aesthetic proportion and symmetry.

  6.  \textbf{Hemostasis and Closure:} Meticulous hemostasis is achieved using electrocautery to prevent hematoma formation. The surgical sites are irrigated. The glandular excision site is closed with absorbable sutures to approximate the deep tissues. The skin incisions are closed with fine absorbable sutures, and often reinforced with sterile strips or skin glue.

  7.  \textbf{Drains and Compression Garment Application:} Small surgical drains (e.g., Jackson-Pratt drains) may be placed temporarily in the surgical bed to collect any fluid or blood, reducing the risk of seroma or hematoma. A compression garment is applied immediately after surgery. This garment minimizes swelling, supports the healing tissues, and helps the skin redrape smoothly over the newly contoured chest.
\end{enumerate}

\section*{Complications \& Risks}
As with any surgical procedure, gynaecomastia surgery carries potential risks, which patients should be thoroughly informed about during the consent process.

\textbf{General Surgical Complications:}
\begin{itemize}
  \item \textit{Anesthesia Risks:} Reactions to anesthetic agents, including nausea, vomiting, allergic reactions, or more serious cardiopulmonary complications such as arrhythmias or respiratory depression.
  \item \textit{Bleeding and Hematoma:} Accumulation of blood under the skin, which may cause pain, swelling, and discoloration, potentially requiring surgical drainage.
  \item \textit{Infection:} Though rare, infection at the surgical site may require antibiotics or, in severe cases, surgical debridement.
  \item \textit{Deep Vein Thrombosis (DVT) and Pulmonary Embolism (PE):} Rare but serious risks associated with any surgery, mitigated by early ambulation and sometimes prophylactic measures like compression stockings.
\end{itemize}

\textbf{Procedure-Specific Complications:}
\begin{itemize}
  \item \textit{Seroma:} Accumulation of clear fluid under the skin, which may require aspiration with a needle or, rarely, drain placement.
  \item \textit{Scarring:} While incisions are typically well-hidden, visible, hypertrophic, or keloid scars can occur, especially in predisposed individuals.
  \item \textit{Nipple-Areola Complex (NAC) Complications:}
    \begin{itemize}
      \item \textit{Nipple Sensation Changes:} Temporary or permanent numbness, hypersensitivity, or partial to complete loss of sensation in the nipples.
      \item \textit{Nipple Necrosis:} Rare, but possible, especially in cases requiring extensive skin removal or free nipple grafts, leading to tissue death.
      \item \textit{Nipple Asymmetry or Malposition:} The nipples may not be perfectly symmetrical in size, shape, or position, or may be inverted.
    \end{itemize}
  \item \textit{Contour Irregularities and Asymmetry:} Unevenness, depressions, persistent fullness, or indentations in the chest contour, potentially requiring revision surgery.
  \item \textit{Skin Laxity:} Especially in older patients or those with significant weight loss, residual skin laxity or wrinkling may be present after tissue removal.
  \item \textit{Fat Necrosis:} Lumps or firmness due to the death of fat cells, which can sometimes be palpable and occasionally require excision.
  \item \textit{Recurrence:} Although rare (typically <5\%), gynaecomastia can recur, particularly if the underlying cause is not addressed, if there is significant weight gain, or if glandular tissue was incompletely removed.
\item \textit{Skin Discoloration:} Bruising and temporary skin discoloration are common, but persistent hyperpigmentation can occur.
\end{itemize}

\section*{Postoperative Care \& Recovery}
Immediately after surgery, patients are transferred to a Post-Anesthesia Care Unit (PACU) for close monitoring as they recover from anesthesia. Pain medication is administered as needed. Most patients are discharged home the same day, though an overnight stay may be recommended for more extensive procedures or specific medical conditions. A compression garment is applied and must be worn continuously for several weeks (typically 4-6 weeks) to reduce swelling, support the healing tissues, and promote optimal skin redraping.

During the first 24-48 hours, patients can expect mild to moderate pain, bruising, and swelling in the chest area. Pain is usually well-controlled with oral analgesics. Drains, if placed, are typically removed within a few days when output significantly decreases. Patients are advised to limit arm movements, avoid lifting heavy objects, and refrain from showering until cleared by the surgeon. Within the first 1-2 weeks, most patients can return to light activities and non-strenuous work. Strenuous exercise, heavy lifting, and activities that put pressure on the chest should be avoided for 4-6 weeks. Bruising usually subsides within 2-3 weeks, while swelling can persist for several months. The final contour will gradually emerge as swelling resolves, with results becoming more apparent at 3-6 months post-surgery. Long-term, patients are advised to maintain a stable weight and healthy lifestyle to preserve their results.

During gynaecomastia surgery, the surgeon meticulously documents the intraoperative findings, including the amount and distribution of glandular versus adipose tissue, the presence of any fibrous bands, and the overall chest contour before and after tissue removal. The excised tissue is then sent for pathological examination.

The pathology report on the resected tissues typically confirms the benign nature of the breast enlargement. Histologically, it often reveals ductal hyperplasia, an increase in the number of breast ducts, accompanied by stromal fibrosis, which is an increase in the fibrous connective tissue. Adipose tissue is also commonly found, reflecting the fatty component of the enlargement. The pathologist will carefully examine the tissue to rule out any atypical hyperplasia or, more importantly, the presence of breast malignancy, which is a rare but critical differential diagnosis. The report provides definitive confirmation of benign gynaecomastia, reassuring both the patient and the surgeon.

A normal and successful postoperative course following gynaecomastia surgery is characterized by a predictable progression of healing and resolution of symptoms. Immediately after surgery, patients will experience mild to moderate pain, managed effectively with prescribed oral analgesics. Swelling and bruising are expected and will gradually subside over several weeks. The compression garment plays a crucial role in minimizing these effects and promoting optimal skin redraping. Within the first few days, any surgical drains are typically removed.

As healing progresses, the chest contour will become progressively flatter and firmer, with the final aesthetic result becoming apparent over 3 to 6 months as all residual swelling resolves. Incisions, initially red and slightly raised, will mature into fine, inconspicuous lines over 6 to 12 months. Nipple sensation, if temporarily altered, usually returns to normal within weeks to months. The expected course includes a gradual return to normal activities, with full strenuous activity typically resumed by 6-8 weeks. The ultimate goal is the complete resolution of preoperative physical symptoms like pain or tenderness, coupled with significant improvement in the patient's psychological well-being and body image, leading to high patient satisfaction.

Postoperative care and follow-up are essential to monitor healing, address concerns, and ensure optimal long-term results after gynaecomastia surgery.

\begin{itemize}
  \item \textbf{Initial Postoperative Visits:} Typically scheduled at 1 week, 2-4 weeks, and 3 months after surgery to check wound healing, remove any drains or non-absorbable sutures, and assess for early complications such as hematoma or seroma.
  \item \textbf{Compression Garment Management:} The surgeon will advise on the duration of compression garment wear, typically 4-6 weeks continuously, and assess its fit and effectiveness.
  \item \textbf{Activity Resumption Guidance:} Gradual return to normal daily activities and exercise will be guided by the surgeon, usually with full activity clearance by 6-8 weeks post-surgery.
  \item \textbf{Scar Management:} Advice on scar massage, silicone sheeting, or other scar management techniques (e.g., topical creams) may be provided after the initial healing phase to optimize scar appearance.
  \item \textbf{Psychological Support:} For patients who experienced significant psychological distress preoperatively, ongoing support or counseling may be beneficial to help adjust to their new body image.
  \item \textbf{Warning Signs:} Patients are instructed to contact their surgeon immediately for signs of infection (e.g., fever, increasing redness, warmth, pus from incisions), excessive or sudden pain, sudden swelling, or persistent numbness that is concerning.
  \item \textbf{Long-term Follow-up:} Annual or less frequent follow-up visits may be recommended to monitor long-term results and address any late concerns.
\item \textbf{Revision Surgery Consideration:} In rare cases of significant contour irregularities, asymmetry, or recurrence, revision surgery may be considered after at least 6-12 months, once all swelling has resolved and tissues have fully settled.
\end{itemize}
Adherence to these postoperative instructions is crucial for achieving the best possible outcome and minimizing complications.

\section*{Outcomes \& Prognosis}
Gynaecomastia is a common condition, affecting a significant portion of the male population across different age groups. It exhibits a trimodal age distribution, with peaks occurring during neonatal life, adolescence, and in older men. Neonatal gynaecomastia is transient, affecting 60-90\% of male infants due to maternal estrogens, and typically resolves within a few weeks. Adolescent gynaecomastia is highly prevalent, affecting 50-70\% of boys during puberty, typically between ages 10 and 16, and often resolves spontaneously within 1-3 years. In adult men, the incidence increases with age, affecting approximately 30-60\% of men over 50, often associated with age-related hormonal changes, obesity, or medication use. The most common indications for surgery are persistent adolescent gynaecomastia and adult idiopathic gynaecomastia. While the condition itself is benign, the psychological morbidity can be substantial, leading to significant distress and impaired quality of life. Surgical mortality is exceedingly rare, primarily associated with general anesthesia risks. Morbidity rates are generally low, with common complications including hematoma (1-5\%), seroma (2-10\%), and temporary nipple sensation changes.

The outcomes of gynaecomastia surgery are generally excellent, with high rates of patient satisfaction and significant improvement in quality of life. Typical success rates, defined by patient satisfaction with aesthetic results and symptom resolution, range from 85\% to over 95\%. Functional outcomes include the ability to wear desired clothing, participate in sports without self-consciousness, and improved posture and body mechanics. Long-term survival outcomes are not directly impacted by the surgery itself, as it addresses a benign condition.

Recurrence rates are low, typically less than 5\%, especially when the underlying cause has been identified and managed, and the patient maintains a stable weight. Prognostic factors for optimal outcomes include appropriate patient selection, meticulous surgical technique, adherence to postoperative care instructions, and a stable body weight. While minor contour irregularities or asymmetry may occasionally necessitate revision surgery, the overall prognosis for achieving a lasting, aesthetically pleasing, and psychologically beneficial result is very favorable. Studies consistently demonstrate that surgical correction of gynaecomastia leads to significant improvements in self-esteem, body image, and overall psychological well-being.

\section*{References}
\begin{enumerate}
  \item MedlinePlus. Gynecomastia Surgery. U.S. National Library of Medicine; 2024. Available from: \url{https://medlineplus.gov/ency/article/002931.htm}
  \item Sabiston DC, Townsend CM. Sabiston Textbook of Surgery: The Biological Basis of Modern Surgical Practice. 21st ed. Elsevier; 2021. (Provides comprehensive surgical principles and specific procedures).
  \item American Society of Plastic Surgeons. Gynecomastia Surgery. Available from: \url{https://www.plasticsurgery.org/cosmetic-procedures/gynecomastia-surgery}
  \item Niewoehner CB, Schorer AE. Gynaecomastia and the Enigma of the Breast. J Clin Endocrinol Metab. 2008;93(3):667-672.
\end{enumerate}

\clearpage
